% Chapter 32: The 2026 Fields Medal — Overview and Outlook
\let\LocalMacro\relax

\chapter{The 2026 Fields Medal: Overview and Outlook}

\section{Overview}
On July 23, 2026, at the International Congress of Mathematicians (ICM) in Philadelphia---the first ICM held in the United States since 1986---the International Mathematical Union announced the four recipients of the 2026 Fields Medal. The medal, often called the ``Nobel Prize of mathematics,'' is awarded every four years to up to four mathematicians under the age of 40.

\begin{historicalbox}
``The four medalists exemplify the depth, originality and vitality of contemporary mathematics,'' said Hiraku Nakajima, president of the IMU, at the ceremony.
\end{historicalbox}

The four medalists---Yu Deng (PDEs), John Pardon (symplectic geometry), Jacob Tsimerman (arithmetic geometry), and Hong Wang (harmonic analysis)---represent a stunning cross-section of modern mathematics. Their work is covered in detail in Chapters 28--31; this chapter places their achievements in broader context.

\section{The Ceremony}

The 2026 ceremony was particularly historic: it marked only the third time a woman received the medal, and it doubled the all-time count of Chinese-born recipients. Additional prizes awarded at the ceremony included:

\begin{itemize}
\item \textbf{Abacus Medal}: Shayan Oveis Gharan (University of Washington) for mathematical contributions to computer science
\item \textbf{Chern Medal}: Graeme Segal (University of Oxford) for lifetime achievement
\item \textbf{Gauss Prize}: Yurii Nesterov (University of Louvain) for applied mathematics
\end{itemize}

\section{The 2025 Mathematical Landscape}

The four 2026 Fields Medalists did not work in isolation. The year 2025 saw a remarkable wave of mathematical breakthroughs across multiple fields.

\subsection{AI and Formal Verification}
The year 2025 marked a turning point for artificial intelligence in mathematics:
\begin{itemize}
\item \textbf{Unit distance conjecture resolved by AI}: An 80-year-old Erd\H{o}s problem---the unit distance conjecture---was resolved using generative AI, which identified a pattern, proved it, and converted the argument into a formally verified proof.
\item \textbf{AlphaProof and DeepSeek-Prover}: AI systems based on reinforcement learning and formal verification (in Lean 4) achieved silver-medal performance at the International Mathematical Olympiad and proved research-level theorems in algebra and number theory.
\item \textbf{Neuromorphic computing for PDEs}: Brain-inspired hardware demonstrated the ability to solve high-dimensional partial differential equations with energy efficiency orders of magnitude better than traditional supercomputers.
\end{itemize}

\subsection{The Geometric Langlands Conjecture}
The geometric Langlands conjecture---one of the deepest unsolved problems in mathematics, connecting number theory, representation theory, and algebraic geometry---saw major progress in 2024--2025. A collaborative effort led by Arinkin, Gaitsgory, and Raskin produced a claimed proof of the unramified case, building on decades of work by Drinfeld, Beilinson, and many others (see Chapter 12).

\subsection{New Mathematical Objects}
\begin{itemize}
\item \textbf{The Noperthedron}: A 90-vertex convex polyhedron discovered by Steininger and Yurkevich that disproves a long-standing conjecture about Rupert's property (the ability to pass a copy of a polyhedron through a hole of the same shape).
\item \textbf{Moving Sofa refinements}: New upper bounds on the moving sofa constant, approaching the theoretical minimum.
\item \textbf{Knot theory advances}: New invariants in knot homology connected to quantum computing.
\end{itemize}

\section{Impact and Outlook}

\begin{itemize}
\item \textbf{Convergence of fields}: The 2026 medals reflect a deepening interconnection between analysis, geometry, number theory, and mathematical physics. Deng's work bridges statistical mechanics and PDE theory; Pardon's connects topology to string physics; Tsimerman's unites model theory with arithmetic geometry; Wang's fuses algebraic geometry with harmonic analysis.
\item \textbf{The role of AI}: As Tsimerman noted in his acceptance remarks, ``There are big changes coming to our world, and I want people to focus on that.'' The integration of AI tools---from formal verification to conjecture generation---is reshaping how mathematical research is conducted.
\item \textbf{Demographic shifts}: The election of two Chinese-born medalists reflects decades of investment in mathematical education in China, producing a new generation of mathematicians at the highest level.
\item \textbf{Open frontiers}: Several of the medalists' techniques open new directions:
  \begin{itemize}
  \item Can Deng's methods for deriving fluid equations be extended to the full range of molecular interactions?
  \item Will Pardon's contact invariants resolve the full Conley--Zehnder conjecture?
  \item Can Tsimerman's approach to Shimura varieties shed light on the Sato--Tate conjecture in higher dimensions?
  \item Will the Kakeya conjecture be settled in dimensions 4 and higher?
  \end{itemize}
\end{itemize}

\section{Exercises}

\begin{exercise}[Basic]
Research the history of the Fields Medal. How many women have won it, and in which years? How does this compare to other major scientific prizes?
\end{exercise}

\begin{exercise}[Intermediate]
Compare the mathematical backgrounds of the four 2026 medalists. What common themes---if any---appear in their approaches to mathematics?
\end{exercise}

\begin{exercise}[Challenge]
The AI breakthroughs of 2025 (unit distance conjecture, AlphaProof) suggest that artificial intelligence can prove research-level theorems. Discuss the potential impact of AI on mathematical research over the next 20 years. Will AI replace mathematicians, or augment them?
\end{exercise}

\section{Timeline}

\begin{tabularx}{\linewidth}{lX}
\toprule
\textbf{Year} & \textbf{Event} \\ \midrule
2024--2025 & Arinkin, Gaitsgory, Raskin et al.\ claim unramified geometric Langlands \\
2025 & AI resolves the unit distance conjecture (Erd\H{o}s, 1946) \\
2025 & AlphaProof achieves silver-medal performance at IMO \\
2026 & ICM held in Philadelphia (first in US since 1986) \\
2026 & Deng, Pardon, Tsimerman, Wang awarded the Fields Medal \\
\bottomrule
\end{tabularx}

\section{Citations}

\begin{enumerate}
\item International Mathematical Union. (2026). ``Fields Medal citations, ICM 2026.'' imu-awards.org.
\item Nature. (2026). ``Rising stars of mathematics awarded prestigious 2026 Fields Medal.'' Nature, d41586-026-02169-1.
\item Scientific American. (2026). ``Four young mathematicians awarded the 2026 Fields Medals.''
\end{enumerate}
